\subsubsection{Cấu trúc dữ liệu \texttt{struct sysinfo}}

Để cung cấp bức tranh toàn diện và chính xác về trạng thái quản lý tập tin và tiến trình trong hệ điều hành xv6, cấu trúc \texttt{struct sysinfo} cùng các cấu trúc con được định nghĩa trong \texttt{kernel/sysinfo.h} theo chuẩn thiết kế đã thống nhất giữa các thành viên (Bảng~\ref{tab:sysinfo}). Process và file được nối với nhau qua chỉ số trong \texttt{ftable}: fd thứ $k$ của một process trỏ tới file có \texttt{idx} bằng \texttt{fidx[}$k$\texttt{]}.

\begin{table}[H]
\centering
\begin{tabularx}{\textwidth}{|l|>{\raggedright\arraybackslash}X|}
\hline
\textbf{Trường} & \textbf{Ý nghĩa} \\ \hline
\multicolumn{2}{|l|}{\textbf{\texttt{struct sysinfo}}} \\ \hline
\texttt{nfile}, \texttt{files[NFILE]} & Số file đang mở trong toàn hệ thống và danh sách của chúng \\ \hline
\texttt{nproc}, \texttt{procs[NPROC]} & Số process có trạng thái khác \texttt{UNUSED} và danh sách của chúng \\ \hline
\multicolumn{2}{|l|}{\textbf{\texttt{struct sysinfo\_file}}: một file đang mở} \\ \hline
\texttt{idx} & Vị trí trong \texttt{ftable.file[]}, dùng để nối với process \\ \hline
\texttt{type} & Loại file: \texttt{FD\_PIPE}, \texttt{FD\_INODE} hoặc \texttt{FD\_DEVICE} \\ \hline
\texttt{readable}, \texttt{writable} & Chế độ mở \\ \hline
\texttt{ref} & Reference count \\ \hline
\texttt{npids}, \texttt{pids[NPROC]} & Các process có ít nhất một fd trỏ tới file này \\ \hline
\multicolumn{2}{|l|}{\textbf{\texttt{struct sysinfo\_proc}}: một process} \\ \hline
\texttt{pid}, \texttt{state} & pid và trạng thái (giá trị của \texttt{enum procstate}) \\ \hline
\texttt{ppid} & pid của cha, bằng 0 nếu không có cha (\texttt{init}) \\ \hline
\texttt{nsched} & Số lần được scheduler chọn chạy \\ \hline
\texttt{name[16]} & Tên process \\ \hline
\texttt{nfd}, \texttt{fds[NOFILE]}, \texttt{fidx[NOFILE]} & Số fd đang mở, các số fd, và \texttt{idx} của file mà mỗi fd trỏ tới \\ \hline
\end{tabularx}
\caption{Các trường của \texttt{struct sysinfo}}
\label{tab:sysinfo}
\end{table}

\subsubsection{Thách thức về Kích thước Ngăn xếp Kernel và Giải pháp \texttt{copyout} từng phần tử}

\textbf{Vấn đề tràn ngăn xếp (Stack Overflow) trong Kernel}:
\begin{itemize}
    \item Kích thước tính toán của các cấu trúc:
    \begin{align*}
    \text{sizeof}(\texttt{struct sysinfo\_file}) &= 4 \times 5 + 4 + 4 \times 64 = 280\text{ bytes} \\
    \text{sizeof}(\texttt{files}[100]) &= 280 \times 100 = 28,000\text{ bytes} \approx 28\text{ KB} \\
    \text{sizeof}(\texttt{struct sysinfo\_proc}) &= 4 \times 4 + 16 + 4 + 16 \times 4 + 16 \times 4 = 164\text{ bytes} \\
    \text{sizeof}(\texttt{procs}[64]) &= 164 \times 64 = 10,496\text{ bytes} \approx 10.5\text{ KB} \\
    \text{Tổng kích thước }\texttt{struct sysinfo} &\approx 28\text{ KB} + 10.5\text{ KB} + 8\text{ bytes} \approx 38.5\text{ KB}
    \end{align*}
    \item Trong xv6, ngăn xếp chế độ nhân (kernel stack) của mỗi tiến trình chỉ được cấp phát đúng \textbf{1 trang bộ nhớ vật lý} duy nhất (\texttt{PGSIZE = 4096 bytes = 4 KB}), và ngăn xếp của tiến trình người dùng cũng chỉ có 1 trang (\texttt{USERSTACK = 1} trong \texttt{kernel/param.h}).
    \item Nếu bất kỳ hàm nào trong kernel khai báo một biến cục bộ dạng \texttt{struct sysinfo info;} trên ngăn xếp, thao tác này sẽ tiêu tốn ngay lập tức $\approx 38.5\text{ KB}$, vượt gấp gần 10 lần dung lượng ngăn xếp cho phép ($4\text{ KB}$). Hậu quả là ngăn xếp lập tức bị tràn qua trang bảo vệ (guard page), làm sập hệ điều hành (kernel panic do page fault).
\end{itemize}

\textbf{Giải pháp Kiến trúc}:
\begin{itemize}
    \item Trong kernel, tuyệt đối \textbf{không} khai báo toàn bộ cấu trúc \texttt{struct sysinfo} trên ngăn xếp.
    \item Thay vào đó, kernel chỉ khai báo \textbf{duy nhất một phần tử} cục bộ (\texttt{struct sysinfo\_file sf} $\approx 280$ bytes trong \texttt{filecollect}, hoặc \texttt{struct sysinfo\_proc sp} $\approx 164$ bytes trong \texttt{proccollect}) trên stack tại mỗi vòng lặp.
    \item Ngay sau khi trích xuất và chuẩn bị xong dữ liệu cho phần tử thứ $k$, kernel sử dụng hàm \texttt{copyout()} để sao chép trực tiếp phần tử đó vào đúng vị trí đích trên không gian địa chỉ ảo của tiến trình người dùng. Với mảng \texttt{files}, vị trí đó là \texttt{addr} cộng offset của trường \texttt{files} trong \texttt{struct sysinfo} (tính bằng \texttt{\_\_builtin\_offsetof}) cộng $k \times$ \texttt{sizeof(struct sysinfo\_file)}; mảng \texttt{procs} tính tương tự.
    \item Ở phía chương trình người dùng (\texttt{sysinfotest.c}), do ngăn xếp user cũng chỉ có 1 trang (4KB), biến \texttt{struct sysinfo} bắt buộc phải được khai báo ở phạm vi \textbf{toàn cục} (được lưu tại phân vùng BSS/DATA) hoặc cấp phát động qua \texttt{malloc()}, không được để làm biến cục bộ trong hàm.
\end{itemize}

\subsubsection{Sơ đồ Luồng Gọi Hệ thống \texttt{sysinfo}}

Quy trình kích hoạt và điều phối của syscall \texttt{sysinfo} được tóm lược như sau:

\begin{enumerate}
    \item \textbf{User Space}: Tiến trình người dùng gọi \texttt{sysinfo(\&info)}.
    \item \textbf{Traps}: Lệnh \texttt{ecall} đưa CPU vào kernel $\rightarrow$ \texttt{usertrap()} $\rightarrow$ \texttt{syscall()} phân phối tới \texttt{sys\_sysinfo()} (\texttt{kernel/sysproc.c}).
    \item \textbf{Điều phối trong \texttt{sys\_sysinfo()}}:
    \begin{itemize}
        \item Gọi \texttt{argaddr(0, \&addr)} để lấy địa chỉ con trỏ buffer người dùng.
        \item Gọi \texttt{nfile = filecollect(addr)} (\texttt{kernel/file.c}): duyệt bảng \texttt{ftable}, thu thập các file đang mở và ghi trực tiếp vào mảng \texttt{info.files} của user.
        \item Gọi \texttt{nproc = proccollect(addr)} (\texttt{kernel/proc.c}): duyệt mảng \texttt{proc}, thu thập các tiến trình đang hoạt động và ghi trực tiếp vào mảng \texttt{info.procs} của user.
        \item Thực hiện \texttt{copyout()} hai biến đếm \texttt{nfile} và \texttt{nproc} vào các trường đầu của cấu trúc tại địa chỉ \texttt{addr}.
        \item Trả về $0$ nếu toàn bộ thành công, hoặc $-1$ nếu có bất kỳ lỗi sao chép bộ nhớ nào (địa chỉ user không hợp lệ).
    \end{itemize}
\end{enumerate}

